\section{Limitations}

While the system successfully implements its core functionalities, several limitations were identified that affect both the user-facing software and the scope of the project itself.
These limitations are discussed in two categories: software limitations and project limitations.

\textbf{System Limitations}

From the perspective of the implemented software, one key limitation is the limited data available. 
The application currently supports only a predefined set of geographic regions and a relatively small number of cuisines. 
Users cannot explore regions or dishes outside of this dataset, which restricts the system's coverage 
and reduces its usefulness for less common areas or cuisines.

Furthermore, the current dataset remains temporary and externally sourced. 
The cuisine information was primarily retrieved and verified through the Wikipedia API, 
rather than being curated, validated, or maintained within a dedicated internal database. 
As a result, the system depends on third-party data availability and consistency, 
which may affect long-term reliability, completeness, and accuracy.

Another limitation lies in the available features. 
Although the system allows users to browse regions and cuisines, submit reviews, and search for restaurants, 
several advanced functionalities are not implemented. 
Features such as region or cuisine search, user personalization, review moderation, ranking mechanisms, 
and real-time updates are missing. 
In addition, the system does not provide visual representations of the cuisines, 
such as images of the dishes, which could significantly enhance user understanding 
and engagement. 
The absence of visual content reduces the overall user experience, 
as food-related applications often rely heavily on imagery to attract interest 
and provide contextual clarity. 
These omissions limit the interactivity and engagement potential of the system.

Finally, while unit and module-level integration testing was conducted, 
comprehensive end-to-end testing (from region selection, cuisine exploration, 
review submission, and restaurant search) was not completed.
This means that potential issues in complex interactions or combined module behavior may remain undetected.


\textbf{Project Limitations}

In terms of the broader project scope, the primary purpose was to explore backend design choices 
related to architecture, database strategies, and caching mechanisms. 
While these topics were researched thoroughly, not all strategies could be implemented in practice. 

Some technical approaches, including distributed system design, alternative caching strategies, 
and different database architectures, were not examined through practical application. 
Although these concepts were researched at a theoretical level, 
they were not implemented within the system. 
If these approaches had been implemented and evaluated in practice, 
they could have provided deeper insight into scalability, performance trade-offs, 
and long-term maintainability. 
However, due to time and scope constraints, practical experimentation with these strategies 
was not applicable within the duration of the project.

These limitations not only define the current scope of the system 
but also establish a foundation for future research and system improvement.

